\subsection*{Grammatical roles transfer more readily than freely worded questions}
We asked whether supervision on the two arguments of a verb helps the writer answer questions about different people. Both training arms saw the same pictured grammatical annotations and active and passive events. Only the correspondence between the question and the pictured answer differed. On the 106 eligible test predicates in Chinese GSD, the matched reader returned the correct subject and object for 36, compared with one for exchanged assignments and zero for the preceding writer (Fig.~\ref{fig:pictured-participant-reader}a). The difference relative to exchanged assignments was 33.0 percentage points (sentence-bootstrap 95\% interval, 24.3--41.9). For predicates absent from the GSD training split, the matched reader recovered both answers in 11 of 36 cases. All natural treebank items retained subject--verb--object order and a fixed pair of grammatical questions.

On 16 new situations with independently worded questions, the matched and exchanged readers each gave 18/64 exact answers. Neither solved all four questions for any situation. On the map-drawing scene, the matched reader identified the person who drew the map. For the person who needed directions it returned the name followed by the clause describing how the person got lost (Fig.~\ref{fig:pictured-participant-reader}c). It retained 148/387 natural person answers, 344/404 familiar compositions and 30/32 descriptions involving bronze names. The number of complete natural person pairs fell from 19 to 16 of 131. All three readers gave 28/504 exact referents on the external Chinese XWinograd split when the options were withheld. The improvement on annotated grammatical arguments did not establish the registered transfer criterion.
